\documentclass[12pt]{article}
\usepackage[T1]{fontenc}
\usepackage{natbib}
\usepackage{setspace}
\usepackage{lineno}
\usepackage{draftwatermark}
\SetWatermarkText{CONFIDENTIAL}
\SetWatermarkScale{0.6}
\doublespacing
\linenumbers
\title{Information Aggregation in Sequential Markets: A Fixture}
\author{Anonymous}
\date{}
\begin{document}
\maketitle
\begin{abstract}
We study how sequential traders aggregate private signals. Our identification strategy exploits the staggered rollout of a disclosure rule. The effect is statistically significant but economically small.
\end{abstract}
\section{Introduction}
Herding has been studied extensively since \citet{banerjee1992} and \citet{bikhchandani1992}. Later work extends the framework to markets \citep{avery1998, smith2000}. As shown by \citet{acemoglu2011}, network structure determines whether learning is efficient; see also \citep[p.~12]{smith2000} and \citet{kogan2020,lee2021}. The coefficient in Table 2 is 0.91 with a confidence interval that crosses one, which is difficult to reconcile with the claim of a robust effect. We define the office-specific identification threshold and afterwards verify the differentiability of the payoff function.

In \citet{banerjee1992}, agents observe only the actions of their predecessors. The \citet{avery1998} model adds a market maker who updates prices. See \citet{lee2021} for heterogeneous priors. Following \citet{smith2000}, we allow unbounded signals. Equilibrium notions are taken from \citet{osborne1994}, and the voting application follows the sequential procedure of \citet{dekel2000}.

\section{Model}
Consider a sequence of agents $i = 1, 2, \ldots$ who each observe a private signal $s_i$ and the history of actions $h_i$. The posterior belief satisfies
\begin{equation}
\mu_i = \frac{\pi \, f(s_i \mid H)}{\pi f(s_i \mid H) + (1-\pi) f(s_i \mid L)} . \label{eq:posterior}
\end{equation}
Proposition 1 states that an information cascade begins with positive probability. The proof follows \citet{bikhchandani1992} with minor modifications. It is easy to see that the inequality is preserved when both sides are multiplied by the likelihood ratio.

\section{Results}
\setcounter{table}{1}
Table~\ref{tab:main} reports the estimates; the posterior in (\ref{eq:posterior}) is evaluated at the sample mean of the signal.
\begin{table}[h]
\centering
\caption{Baseline estimates}\label{tab:main}
\begin{tabular}{lcc}
\hline
Specification & Estimate & 95\% interval \\
\hline
Baseline & 0.91 & [0.78, 1.06] \\
Staggered rollout & 0.87 & [0.74, 1.02] \\
Office fixed effects & 0.95 & [0.80, 1.12] \\
\hline
\end{tabular}
\end{table}

\bibliographystyle{plainnat}
\begin{thebibliography}{9}
\bibitem[Acemoglu et~al.(2011)Acemoglu, Dahleh, Lobel, and Ozdaglar]{acemoglu2011}
Acemoglu, D., Dahleh, M.~A., Lobel, I., and Ozdaglar, A. (2011).
\newblock Bayesian learning in social networks.
\newblock {\em Review of Economic Studies}, 78(4):1201--1236.
\bibitem[Avery and Zemsky(1998)]{avery1998}
Avery, C. and Zemsky, P. (1998).
\newblock Multidimensional uncertainty and herd behavior in financial markets.
\newblock {\em American Economic Review}, 88(4):724--748.
\bibitem[Banerjee(1992)]{banerjee1992}
Banerjee, A.~V. (1992).
\newblock A simple model of herd behavior.
\newblock {\em Quarterly Journal of Economics}, 107(3):797--817.
\bibitem[Bikhchandani et~al.(1992)Bikhchandani, Hirshleifer, and Welch]{bikhchandani1992}
Bikhchandani, S., Hirshleifer, D., and Welch, I. (1992).
\newblock A theory of fads, fashion, custom, and cultural change as informational cascades.
\newblock {\em Journal of Political Economy}, 100(5):992--1026.
\bibitem[Dekel and Piccione(2000)]{dekel2000}
Dekel E, Piccione M (2000) Sequential voting procedures in symmetric binary elections.\\
J Polit Econ 108(1):34--55
\bibitem[Kogan(2020)]{kogan2020}
Kogan, D. (2020).
\newblock Learning from reviews with strategic sellers.
\newblock Working paper.
\bibitem[Lee(2021)]{lee2021}
Lee, J. (2021).
\newblock Observational learning with heterogeneous priors.
\newblock {\em Games and Economic Behavior}, 125:1--20.
\bibitem[Osborne and Rubinstein(1994)]{osborne1994}
Osborne, M.~J. and Rubinstein, A. (1994).\\
\newblock {\em A Course in Game Theory}. MIT Press.
\bibitem[Smith and S{\o}rensen(2000)]{smith2000}
Smith, L. and S{\o}rensen, P. (2000).
\newblock Pathological outcomes of observational learning.
\newblock {\em Econometrica}, 68(2):371--398.
\end{thebibliography}
\end{document}
